% =====================================================================
% Chapter 8 -- Results
% =====================================================================
\documentclass[../preamble.tex]{subfiles}
\begin{document}

\chapter{Results}
\label{ch:results}

\section{Headline numbers}
\label{sec:results-headline}

The pre-HPO ensemble is the recommended deployed model. The headline metrics
on the held-out PCOSgen test set are:

\begin{table}[htbp]
  \centering
  \small
  \caption{Pre-HPO ensemble headline metrics on the held-out PCOSgen test set ($n=1468$).}
  \label{tab:headline}
  \begin{tabular}{lr}
    \toprule
    \textbf{Metric} & \textbf{Value} \\
    \midrule
    AUC-ROC (95\% CI) & 0.9487 (0.9368--0.9596) \\
    F1 & 0.8665 \\
    MCC & 0.7884 \\
    Brier & 0.0981 \\
    ECE (before pass 1) & 0.0854 \\
    ECE (after pass 2) & 0.0810 \\
    \bottomrule
  \end{tabular}
\end{table}


The 95\% confidence interval on AUC is $[0.9368, 0.9596]$, computed by
bootstrap resampling. F1 and MCC reflect the same operational balance: high
sensitivity, moderate specificity, and a Brier score consistent with
calibrated probabilities.

\section{Stage-wise confusion matrices}
\label{sec:results-confusion}

\Cref{fig:confusion-stages} shows the three-stage confusion-matrix
progression: foundation zero-shot, fine-tuned best single model, and
pre-HPO ensemble. The numerics are in \cref{tab:stage-cm}.

\begin{figure}[htbp]
  \centering
  \includegraphics[width=\linewidth]{figures/paper_fig_confusion_stages.png}
  \caption{Three-stage confusion-matrix progression. Stage 1 was scored on
           the full PCOSgen corpus ($n=4668$); Stages 2 and 3 share the
           held-out test set ($n=1468$).}
  \label{fig:confusion-stages}
\end{figure}

\begin{table}[htbp]
  \centering
  \small
  \caption{Confusion-matrix counts per pipeline stage across evaluations.}
  \label{tab:stage-cm}
  \begin{tabular}{lrrrrr}
    \toprule
    \textbf{Stage} & \textbf{TN} & \textbf{FP} & \textbf{FN} & \textbf{TP} & \textbf{$n$} \\
    \midrule
    1. Foundation (SRAD + ResNet-50) & 832 & 209 & 1245 & 2382 & 4668 \\
    2. Fine-tune (SRAD + DenseNet-121) & 812 & 107 & 60 & 489 & 1468 \\
    3. Pre-HPO ensemble & 838 & 81 & 48 & 501 & 1468 \\
    \bottomrule
  \end{tabular}
\end{table}


The headline observation is that the foundation stage fails in the
screening-relevant direction: it predicts PCOS for almost every image
because the Figshare training distribution is 3.92:1 PCOS-heavy. After
fine-tuning on the PCOSgen train pool, the false-positive count drops from
998 to 241 (DenseNet-121 alone) and to 164 (ensemble). The ensemble's
specificity improves from 0.041 to 0.822 while keeping sensitivity above
0.99.

\section{Internal and external AUC curves}
\label{sec:results-internal-external}

\Cref{fig:internal-pad} shows the per-architecture internal AUC on the
Figshare test split. All eighteen configurations plateau in the 0.99 range.
This saturation is the central evidence that internal accuracy alone cannot
characterise the pipeline.

\begin{figure}[htbp]
  \centering
  \includegraphics[width=\linewidth]{figures/f1_internal_auc.png}
  \caption{Internal Figshare test-set AUC across all 18 foundation-stage
           configurations. All architectures saturate in the 0.99 range.}
  \label{fig:internal-pad}
\end{figure}

\Cref{fig:external-collapse} shows the corresponding zero-shot AUC on the
PCOSgen corpus: every foundation checkpoint collapses to near-random
performance. The collapse is symmetric across both denoisers.

\section{Pre-HPO versus post-HPO ensemble}
\label{sec:results-pre-post-hpo}

The fine-tune stage produced two ensembles: one using the pre-HPO checkpoints
and one using the post-HPO retrain. The pre-HPO ensemble is the recommended
model because Optuna's small validation set destabilised the retrain for
some architectures; the differences are recorded in the
\texttt{ensemble/top3\_pre\_hpo/} and \texttt{ensemble/top3\_hpo/} artefact
directories.

\begin{figure}[htbp]
  \centering
  \includegraphics[width=\linewidth]{figures/paper_fig_ensemble_metrics.png}
  \caption{Pre-HPO versus post-HPO ensemble on the held-out test set.
           Pre-HPO wins on every metric; this is the recommended model.}
  \label{fig:ensemble-metrics}
\end{figure}

\section{Generalisation recovery}
\label{fig:recovery}
\label{sec:results-recovery}

The generalisation-recovery figure traces the AUC arc from foundation to
fine-tune to ensemble, showing that the target-cohort fine-tune is what
restores the discriminating signal, not the choice of backbone or denoiser.

\begin{figure}[htbp]
  \centering
  \includegraphics[width=\linewidth]{figures/paper_fig_generalization_recovery.png}
  \caption{Generalisation recovery arc: 0.45--0.59 (foundation zero-shot)
           to 0.52--0.77 (no-pad fine-tune) to 0.93 (best single fine-tuned)
           to 0.9487 (pre-HPO ensemble).}
  \label{fig:recovery}
\end{figure}

\section{Operating points}
\label{sec:results-operating}

The operating-point comparison in \cref{fig:sens-vs-fpr} shows that the
default 0.5 threshold has specificity between 0.20 and 0.45 across the 18
configurations, which is too low for a screening tool. The recomputed
operating points for the pre-HPO ensemble are:

\begin{itemize}
  \item default 0.5: sensitivity 0.9927, specificity 0.8215;
  \item Youden-J (threshold 0.6198): sensitivity 0.9909, specificity 0.8324;
  \item sensitivity-targeted (threshold 0.7597): sensitivity 0.9500,
        specificity 0.8553.
\end{itemize}

\begin{figure}[htbp]
  \centering
  \includegraphics[width=\linewidth]{figures/f9_sens_vs_fpr.png}
  \caption{Sensitivity versus false-positive rate at the default 0.5
           threshold. Most configurations cluster near sensitivity 1.0 but
           at the cost of specificity in 0.20--0.45.}
  \label{fig:sens-vs-fpr}
\end{figure}

\section{ROC, PR, and combined curves}
\label{sec:results-curves}

\begin{figure}[htbp]
  \centering
  \includegraphics[width=\linewidth]{figures/paper_fig_combined_roc.png}
  \caption{ROC curves of the three ensemble members and the pre-HPO
           ensemble. The ensemble curve encloses the largest area.}
  \label{fig:combined-roc}
\end{figure}

\begin{figure}[htbp]
  \centering
  \includegraphics[width=\linewidth]{figures/paper_fig_combined_pr.png}
  \caption{Precision-recall curves. The pre-HPO ensemble reaches an
           average precision of 0.8702.}
  \label{fig:combined-pr}
\end{figure}

\section{Per-class precision, recall, F1}
\label{sec:results-perclass}

\begin{figure}[htbp]
  \centering
  \includegraphics[width=\linewidth]{figures/thesis_per_class_pr.png}
  \caption{Per-class precision, recall, and F1 on the held-out test set at
           the default 0.5 threshold. The ensemble achieves the highest
           precision (0.769) and F1 (0.867) of the four configurations.}
  \label{fig:per-class}
\end{figure}

The precision uplift at the ensemble level is the cleanest effect: each
single model produces a precision below 0.35, while the ensemble's
precision rises to 0.77 because the probability average makes the boundary
between PCOS and healthy predictions sharper.

\section{Sweep matrix}
\label{sec:results-sweep}

The 18-run sweep matrix records per-arch-per-prep external AUC, MCC, F1,
specificity, and Brier scores.

\begin{figure}[htbp]
  \centering
  \includegraphics[width=\linewidth]{figures/paper_fig_sweep_matrix.png}
  \caption{Sweep matrix heatmap. The foundation stage saturates internally
           and collapses externally.}
  \label{fig:sweep-matrix}
\end{figure}

\section{Uncertainty summary}
\label{sec:results-unc}

The MC-Dropout entropy summary is in \cref{tab:mc-entropy}. The mean
predictive entropy over the held-out test set is 0.234 nats for DenseNet-121,
0.215 nats for ConvNeXt-Tiny, and 0.160 nats for ViT-B/16. The median
entropies are 0.151, 0.139, and 0.089 respectively.

\begin{table}[htbp]
  \centering
  \small
  \caption{MC-Dropout predictive-entropy summary on the held-out test set.}
  \label{tab:mc-entropy}
  \begin{tabular}{lrr}
    \toprule
    \textbf{Architecture} & \textbf{Mean entropy (nats)} & \textbf{Median entropy (nats)} \\
    \midrule
    DenseNet-121   & 0.234 & 0.151 \\
    ConvNeXt-Tiny  & 0.215 & 0.139 \\
    ViT-B/16       & 0.160 & 0.089 \\
    \bottomrule
  \end{tabular}
\end{table}

\section{Calibration results}
\label{sec:results-calib}

Per-model temperatures and the ensemble temperature are recorded in
\cref{tab:calibration-detail} of \cref{ch:trust}. The reduction in
Expected Calibration Error from 0.0854 to 0.0810 corresponds to a
re-alignment of the reliability diagram. The reliability figures are
reproduced from \cref{fig:calibration}.

\end{document}